\section{The factor two in the finite character moment}\label{sz:shells}
The projection to good types in Section~\ref{B:arithmetic} uses the
unnormalized sum $\sum_{g\in G}|\psi_U^\natural(g)|^2\le2$. The small
constant survives integration over the product of the chosen places and
their conjugates. We give the passage from the estimates for the frame
kernel to this finite sum, with the assumptions on that kernel made explicit.

Let $G_{\rm loc}$ be the unitary group over the semilocal algebra and $K$ its integral subgroup, $K_M$ its
principal congruence subgroup, and $\mu$ the Haar measure of the filling
datum. The character $\chi_M$ is extended by zero off $U K_M$, where
$U$ is the full Bessel unipotent subgroup of $G_{\rm loc}$. On $K$ the support restricts to $(U\cap K)K_M$. It is constant on
right $K_M$-cosets, has modulus at most one, and equals one on $K_M$.
For $K_M\subset Z\subset K$ a compact open intermediate group, set
\[
 a_Z(x)=\int_Z\chi_M(yxy^{-1})\,d\mu(y),\quad
 R_Z=\frac{\mu(Z)\mu(K_M)}{u_F^2},\quad h_F=\varrho+1,
\]
where $u_F>0$ is the unipotent mass of the filling and $h_F$ is its frame
parameter. On $K_M$ the integral equals $\mu(Z)$ by normality.

Here are the precise kernel estimates used on the complement.
In the gauge of scale one, let $P_Z(x)$ be the nonnegative majorant
of the zero-floor frame correlation with empty heads, and let $W_Z(x)$
be its weighted square product. These are the functions
\texttt{formMajorant} and \texttt{weightedSquareProduct} of the filling.
The following four geometric inputs specify the remaining part of F5.
\begin{enumerate}
\item The pairing-kernel estimate gives $|a_Z(x)|^2\le u_F^2 P_Z(x)$.
\item At an active point $y$ and each frame index $i$, the set of stable
products has volume at most the modulus of some element $c$ in it with
$c$ integral and invertible in the semilocal algebra. If every such set
contains a $c$ with integral inverse, the frame reconstruction makes
$(x-1)/M$ an integral matrix. Here the comparison uses the scale-one gauge,
where every bound on matrix size is integrality.
\item If every active point at $x$ has some index with stable-product
volume at most $e<1$, the zero-floor majorant satisfies
$P_Z(x)^2\le e^{2s_F}W_Z(x)$, where $s_F$ is the volume exponent of
the filling.
\item For a measurable $0\le F\le P_Z$ obeying
$F(x)^2\le e^{2s_F}W_Z(x)$ on the positions supporting the majorant,
the dyadic denominator-shell estimate is
\[
 \int F\,d\mu\le e^v R_Z\,2^{10^9h_F^4}
       \qquad(v+10^8h_F^4\le k,\;v\in\mathbb N).
\]
Here the shells group the denominator ideals of the entries of $(x-1)/M$
by the ceiling of the base-two logarithm of their index. The rank $k$
in this assertion is the rank in the filling shape.
\end{enumerate}
These inputs are required for the specific large lattice and its normalized
measure. The following calculation derives the complementary integral from
them and shows which rank margin is spent absorbing its loss.

Suppose $Q\ge2$ and each integral $c$ invertible in the algebra
whose inverse is nonintegral has modulus at most $Q^{-1}$. One can
always take $Q=2$ here. Indeed multiplication by $c$ takes the
integer ring $R$ to the proper finite-index ideal $cR$, and
its modulus equals $[R:cR]^{-1}$. If $x\notin K_M$ and every stable
product set at some active point had volume greater than $Q^{-1}$,
the witnesses $c$ in the second input would all have integral inverse.
Then $(x-1)/M$ would be integral, which in the unimodular integral
model puts $x$ in $K_M$. Thus the third input applies on the complement.
Put $F=\mathbf1_{K_M^c}P_Z$, $B=10^9h_F^4$ and $v=t+B$.
Under $t+2B\le k$ the last input gives
\[
 \int_{K_M^c}P_Z\,d\mu\le Q^{-t}(2/Q)^B R_Z\le Q^{-t}R_Z .
\]
Using the first input there and the exact integrand on $K_M$ yields
\[
 \int |a_Z(x)|^2\,d\mu(x)
 \le\mu(Z)^2\mu(K_M)+Q^{-t}R_Zu_F^2
 =\mu(Z)\mu(K_M)(\mu(Z)+Q^{-t}).
\]

Take $Z=K$, $Q=2$ and $t=h_F$. The filling rank condition is
$h_F+2\cdot10^9 h_F^4\le k$, and the local volume estimate for the
chosen places is $\mu(K)\ge2^{-h_F}$. Hence the parenthesis is at most
$2\mu(K)$. Choose lifts $s_g$ of the finite reduction elements. The
cosets $s_gK_M$ are disjoint. On each, $a_K$ is constant and equals
$\mu(K)\psi_U^\natural(g)$: integration of the conjugated character
is the uniform conjugacy average on the finite quotient. Each coset
has the same positive volume. Summing their integrals and dividing by
$\mu(K)^2\mu(K_M)$ proves \eqref{B:charmean}. Using the whole
product at once gives a single factor two.

\subsection{Integral recovery of a position from its frame}
\label{app:local-frame-recovery}
We prove the stable-product and reconstruction assertions in the second
input of Appendix~\ref{sz:shells}.  The argument takes place over the
entire product of chosen finite places.  Let $A=\prod_v A_v$ be their
product of local fields, $R=\prod_v\mathcal O_v$, and normalize additive
measure by $\operatorname{vol}(R)=1$.  Write $L$ for the coordinate
$R$-lattice of the local Hermitian space and $h$ for its form, linear in
the second argument.  The form is integral and unimodular on $L$.
Write $E_i=\sum_{a<i} A e_a$ for its isotropic frame flag, with
$0\le i\le\varrho$, and $Y_i=E_i^\perp/E_i$.  The Bessel unipotent
group $U$ preserves the flag, fixes $e_i$ modulo $E_i$, and acts
as the identity on $Y_{\varrho}$.  These are the diagonal blocks
of the defining upper unitriangular matrices.  Let $K$ be the
stabilizer of $L$, and $K_M$ its kernel on $L/ML$.  The level $M$
is integral, fixed by the involution and nonzero in each factor.

For $g,g'=gx$ and $w\in E_i^\perp$, call $w'$ a matching lift if
\[
  w'-w\in E_i,\qquad g^{-1}w\in L,\qquad
  g'^{-1}w'-g^{-1}w\in ML .
\]
Define $I_i$ to be the set of values $h(e_i,w)$ for $w$ admitting such
a lift.  At an active point there exist $u,v\in U$ such that
\[
 C=ug\in K,\qquad C'=vg'\in K,\qquad C^{-1}C'\in K_M .
 \tag{*}
\]
Indeed the finite Bessel vector at $y$ has coordinates
$\chi_M(yc^{-1})$ indexed by $c\in Z/K_M$. A nonzero inner product at
$g,g'$ gives one coordinate in which both entries are nonzero. Its
support is $U K_M$, so one can choose $u,v$ with
$ugc^{-1},vg'c^{-1}\in K_M$; now use $c\in K$ and normality.

\begin{proposition}
At each active point $I_i=c_iR$ for some $c_i\in R\cap A^\times$,
and its depth is $[R:I_i]$. If every $I_i$ has index one, then
$x\in K_M$.
\end{proposition}
\begin{proof}
The pairs $(w,w')$ with $w\in E_i^\perp$ and $w'-w\in E_i$
form an $A$-subspace. Intersect it with the inverse image of
$L\oplus L$ under the invertible linear map
\[
 (w,w')\longmapsto (g^{-1}w,M^{-1}(g'^{-1}w'-g^{-1}w)).
\]
The intersection is a full lattice in this subspace: over each DVR
one clears denominators of a basis, and a submodule of a finite free
module is finitely generated. Its image by $(w,w')\mapsto h(e_i,w)$
is a full fractional ideal in every factor, since $(e_i^*,e_i^*)$
has image one for the hyperbolic dual $e_i^*$. Thus $I_i=c_iR$
for a componentwise nonzero generator. Put $f=u^{-1}e_i$ and
$f'=v^{-1}e_i$. These represent $e_i$ modulo $E_i$, and
$g^{-1}f=C^{-1}e_i$, $g'^{-1}f'=C'^{-1}e_i$ both belong to $L$.
For $w\perp E_i$ with $g^{-1}w\in L$,
$h(e_i,w)=h(g^{-1}f,g^{-1}w)\in R$. This gives $I_i\subset R$.
Its inverse volume, the depth by definition, is its index.

Suppose now $I_i=R$ for all $i$. We prove by decreasing induction on
$i$ that every $w\in E_i^\perp$ with $g^{-1}w\in L$ has a matching
lift. For $i=\varrho$ take $w'=v^{-1}u w$. Its class agrees with $w$
by the action on $Y_{\varrho}$. Moreover $g'^{-1}w'=C'^{-1}C g^{-1}w$,
which is congruent to $g^{-1}w$ by $(*)$.

For the induction step choose $z\in E_i^\perp$ and its matching
lift $z'$ with $h(e_i,z)=1$. Put $\alpha=h(e_i,w)\in R$ and
$w_0=w-\alpha z$. Then $w_0\perp E_{i+1}$ and $g^{-1}w_0\in L$.
By induction it has a matching lift $w_0'$ modulo $E_{i+1}$.
Write $w_0'-w_0=\mu e_i+t$ with $t\in E_i$. Then
\[
 \mu=h(z',w_0')-h(z,w_0)\in MR.
\]
Indeed $z'-z\in E_i$, $w_0'\perp E_i$, and $h(z,e_i)=1$.
The two pairings are congruent modulo $M$ on applying $g'^{-1}$ and
$g^{-1}$ respectively, since both pairs of vectors are integral and
congruent modulo $ML$. Now $w'=\alpha z'+w_0'-\mu f'$ has the same
class as $w$ modulo $E_i$, and the subtracted term is zero modulo
$ML$ after applying $g'^{-1}$. This proves the induction.
At $i=0$ the only lift of $w$ is $w$ itself. Taking $w$ in the
columns of $g$ gives $(x^{-1}-1)L\subset ML$. Taking the Hermitian
adjoint and using unimodularity gives $(x-1)L\subset ML$ as well.
Thus $x\in K_M$.
\end{proof}

In particular every active point at $x\notin K_M$ has some $i$ with
$\operatorname{vol}(I_i)\le 1/2$. More precisely the volume is
$\prod_v q_v^{-v(c_{i,v})}$, where $q_v$ is the residue cardinality,
so $1/2$ can be replaced here by $(\min_v q_v)^{-1}$.
This establishes the second geometric input for the stated finite
unimodular model. Converting this depth into a high power of the
volume saving still requires the correlation and integration estimates
in the remaining inputs.

\paragraph{The zero-floor pairing and its normalization.}
Let $\eta$ be Haar measure on $U$, let $u_F=\eta(U\cap K_M)$ and
let $d\dot y$ be quotient measure on $U\backslash G_{\rm loc}$,
normalized by $d\mu=d\eta\,d\dot y$. Write
$V_c(y)=\chi_M(yc^{-1})$ for $K_Mc\in K_M\backslash Z$, and
$H_x(y)=\sum_c\overline{V_c(y)}V_c(yx)$. Then
\[
 a_Z(x)=u_F\int_{U\backslash G_{\rm loc}} H_x(y)\,d\dot y .
\]
In fact, if $u y\in K_Mc$, the two coordinates at $c$ have the
common phase $\chi_M(u^{-1})$. Their product after conjugating the
first equals $\chi_M((uy)x(uy)^{-1})$ by right $K_M$-invariance.
For fixed $y,c$ the set of such $u$ is empty or one right coset
of $U\cap K_M$ within $U$. Its measure in the latter case is $u_F$,
as $U$ is unimodular. If the set is empty then $V_c(y)=0$.
Partitioning $\{u:uy\in Z\}$ by these cosets and applying the
quotient integral formula gives the identity. All coordinates are
bounded with compact support modulo $U$.
In the frame integration coordinates of the filling let $d\theta_0$
be the first floor measure and $d\theta_{>0}$ the measure of the
remaining floors, so $\dot y$ is represented by $y'\sigma$ with
measure $d\theta_{>0}(y')d\theta_0(\sigma)$ on $U$-invariant
integrands. At zero floor the two heads coincide, and the definition
of the majorant therefore reads
\[
  P_Z(x)=
 \left(\int\left|\int H_x(y'\sigma)\,d\theta_{>0}(y')\right|
                  \,d\theta_0(\sigma)\right)^2 .
\]
The triangle inequality now gives $|a_Z(x)|^2\le u_F^2 P_Z(x)$,
the first input of Appendix~\ref{sz:shells}, with coefficient one
for these normalized floor measures.

\paragraph{Counting the denominator ideals.}
Write $n=\varrho+1$ and $H=2^j$ for a dyadic height, $j\ge0$.
The denominator ideal of a matrix $D$ over $A$ is
$J=\{c\in R:cD\text{ is integral}\}$. It is generated by an element
$d\in R\cap A^\times$, componentwise by a uniformizer to the
largest of zero and the pole orders of the matrix entries.
Take $D=(x-1)/M$. Every $x$ belongs to the shell of the least dyadic
$H\ge [R:J]$ and the scale $\bar d d$ makes $(\bar d d)D$ integral.
The scale is invariant under the involution, and its additive modulus
is $[R:J]^{-2}$, since the involution preserves the measure of $R$
and carries ideals to ideals of the same index. Hence this region
costs at most $2 H^2$ in the gauge $2/|\bar d d|$.

There are at most $2H$ such ideals for the chosen places and all
sufficiently large cutoffs. In fact there are at most $4f$ local
factors, since they lie above the two rational primes $p_0,p_1$
and $[E:\mathbb Q]=2f$; any chosen subproduct obeys the same bound.
Each $q_v\ge b^{10}$. Therefore
\[
 \sum_{0<[R:J]<\infty}[R:J]^{-1}
   =\prod_v(1-q_v^{-1})^{-1}
   \le \exp(8f b^{-10})\le 2.
\]
The logarithmic bound uses $-\log(1-s)\le2s$ for $0\le s\le 1/2$.
The last inequality follows from $f\le q_1\le x^{12}$ eventually
and the scales of Theorem~3.1. Every ideal of index at most $H$
contributes at least $1/H$. This gives the required covering count,
indeed with $2H$ in place of $2H^{9n}$.

\paragraph{The finite quotient only needs positions in the integral group.}
The moment in (10.11) is a sum over $K/K_M$, so restrict the outer
integral to $x\in K$. In this region a stronger form of flag recovery
gives the decay needed for summing over $J$. Write
$J(x)=\{c\in R:c(x-1)/M\text{ is integral}\}$ and
$N(x)=[R:J(x)]$. Each chosen place is unramified over its level
prime: the two level primes are prime to the conductor $q_1d$ of a
cyclotomic field containing $E$. Thus $M=p_0p_1$ has valuation one
in each local factor. Since $x$ is integral, the local denominator
ideal is either the whole valuation ring or its maximal ideal. The
latter occurs precisely when the reduction of $x$ at that component
differs from the identity.

Decompose $R$ according to the orbits of its factors under the
involution. For each orbit $o$ contributing to $N(x)$, write $q_o$
for the residue cardinality of any member. This is independent of the
member. The orbit contributes at most $q_o^2$ to $N(x)$.
At every active pair the hypotheses $(*)$ project to each orbit.
The matching conditions defining $I_i$ are componentwise over these
orbits, so $I_i$ is the product of the corresponding ideals $I_{i,o}$.
If $x$ differs from the identity modulo $M$ on $o$, the preceding
proposition applied there gives an $i$ with
$[R_o:I_{i,o}]\ge q_o$. Multiplication over contributing orbits gives
\[
  \prod_{i<\varrho}[R:I_i]\ge\prod_{o\text{ contributing}}q_o
       \ge N(x)^{1/2}.
\]
Hence, for $\varrho>0$, some index has depth at least
$N(x)^{1/(2\varrho)}$. For $\varrho=0$, the induction at the top
quotient already implies $x\in K_M$ wherever the pairing is active.

For precision we isolate the remaining analytic estimate on one shell.
For $H=2^j$ let $S_H$ be the denominator shell introduced above.
The normalized one-frame bound used here is that if $F$ is
nonnegative, is supported in $S_H$, and satisfies $F^2\le t^2 W_Z$
there, then
\[
 \int F\,d\mu\ \le\
    t\,R_Z\,2^{2\cdot10^6 n^3+1287n} H^{10^5 n^2}.
 \tag{**}
\]
It suffices to impose this statement on functions dominated by $P_Z$
and supported in the positions of the frame correlation.
This inequality is an analytic input. So is the third geometric
input in Appendix~\ref{sz:shells}, giving the weighted pointwise
gain $P_Z^2\le e^{2s_F} W_Z$ if every active frame has a depth at least
$e^{-1}$. No estimate for deep vectors at positions outside $K$ is
required below.

Take $Z=K$ and $F=\mathbf 1_{K\cap S_H}P_K$ for $H\ge2$. On this
shell $N(x)\ge\sqrt H$: for $H=2$ this follows from integrality of
$N(x)>1$, and for $H\ge4$ it follows from $N(x)>H/2$.
The previous product estimate permits $e=H^{-1/(4\varrho)}$.
Consequently $(**)$, with $t=e^{s_F}$, gives
\[
 \int_{K\setminus K_M}P_K\,d\mu\le
 R_K\,2^{2\cdot10^6n^3+1287n}
    \sum_{j\ge1}2^{-j(s_F/(4\varrho)-10^5 n^2)} .
\]
The geometric series pays for its prefactor. In fact, if the block
parameter of the large lattice datum is $m_0+1$, the starting
coordinate rank is $k-2m_0-3$, two blocks each of rank $2m_0+4$
are removed, and the exponent subtracts a further $3n+m_0+2$.
Thus $s_F=k-7m_0-3n-13$, with $m_0<9n$. Since
$k\ge n+2\cdot10^9 n^4$, this gives $s_F\ge10^9 n^4$.
Write $A_0=s_F/(4\varrho)-10^5 n^2$.
Using $\varrho<n$ gives
$A_0\ge 2\cdot10^6n^3+1288n+1$.
The sum is at most $2^{1-A_0}$, so the displayed integral is at most
$2^{-n}R_K$. The zero-floor pairing and the exact value on $K_M$
now give
\[
 \int_K |a_K|^2\,d\mu\le
  \mu(K)\mu(K_M)\bigl(\mu(K)+2^{-n}\bigr).
\]
With the normalized volume bound $\mu(K)\ge2^{-n}$ from
Appendix~\ref{sz:shells}, division by $\mu(K)^2\mu(K_M)$ gives (10.11).
This proves the finite-quotient reduction from the weighted pointwise
gain, the one-shell integration bound $(**)$ and the specified
measure normalizations. Other uses of an integral over the entire
local group require their own estimates outside $K$.

